\section{从早期统一模型到响应与波函数}\label{sec:1-3}

同一个原子核，为什么既能显出单粒子壳层的规则，又能表现出整体转动与振动的特征？上一节区分了实验室电荷关联与内禀形状，但给形状选好坐标，还没有回答它怎样运动。早期统一模型把这个问题变成了可计算的粒子---集体耦合；能量与电磁跃迁的对照，又使其近似接受检验。沿着这条问题线，可以理解后来为何要进一步计算波函数对扰动的响应。

\subsection{核场为什么也需要成为动力学对象}\label{sec:1-3-1}

独立粒子近似把复杂相互作用先归入平均势，核子占据其中的单粒子轨道。Bohr 与 Mottelson 在 1953 年长篇论文的开头，以壳层规则、四极矩和电四极跃迁等现象提出了一个联系：束缚核子的核场由核子自身产生，核子密度的整体变化也会改变单粒子运动。因此，需要研究一个能够变形的壳结构，而不能始终把平均势当成固定背景~\cite{bohr-mottelson-1953-collective-individual-original}。

这里应区分独立粒子近似与相互作用壳模型。后者保留残余相互作用，允许不同占据组态相干混合，在适当模型空间中能够产生集体关联。引入集体坐标，是把许多核子共同参与的变化组织成便于处理的变量，并不意味着壳模型原则上无法描述集体性。两种语言需要联系，因为它们描述的是同一批核子的运动。

\subsection{从形状坐标到有条件的慢运动}\label{sec:1-3-2}

Bohr 的 1952 年论文从小幅表面形变出发，用 $\beta$、$\gamma$ 描写四极形状的幅度和非轴对称程度，用欧拉角描写主轴的空间取向。这是包含旋转变换的非线性坐标变换。在其模型中，动能可以写成形状变化与转动两部分，但转动惯量仍依赖形状，量子波函数还需满足主轴重标记的一致性条件。坐标按形状和取向组织起来，不等于波函数或动力学已经解耦~\cite{bohr-1952-surface-particle-coupling}。

为了给动能提供具体系数，Bohr 采用了不可压、常密度、无旋流的基准。无旋流指所假定速度场的旋度为零，并非物质像刚体一样随表面整体转动。这个流动假设给出振动质量参数与转动惯量，使能谱可以计算；它也增加了需要检验的物理条件。原文已经提醒，核的可形变性会受壳结构影响，流体估计是简化模型。

1953 年的进一步论述把集体坐标明确写为核子坐标的对称函数，并说明绝热近似的能标条件：当粒子激发比所讨论的集体运动快得多时，可以近似认为粒子结构跟随缓慢变化的形状。遇到容易激发、能标可比的粒子自由度，就要显式保留它们与集体运动的耦合。强耦合解也有涨落和混合的适用范围。内禀图像由这些条件获得计算意义，不能读成实验室精确态的一幅固定表面图~\cite{bohr-mottelson-1953-collective-individual-original}。

\subsection{同一个样本，为什么得到不同的形变}\label{sec:1-3-3}

统一模型的解释力也体现在它使不同观测能够相互核对。Bohr 与 Mottelson 在 1953 年论文的 Table~XXIII 中列出 $^{160}\mathrm{Dy}$ 等历史样本：一条路径由转动能级间距约束惯量，再借助流体质量参数映射到形变；另一条路径由寿命所约束的电四极（$E2$）跃迁几率，借助内禀四极矩与几何形变的关系作估计。$^{160}\mathrm{Dy}$ 行中，由能量路径得到的形变高于 $E2$ 路径的结果。两条路径都包含模型映射，表内采用量也属于当时的资料~\cite{bohr-mottelson-1953-collective-individual-original}。

作者在紧邻讨论中既指出四极矩与形变的简单几何关系可能不足，也考虑了质量参数被低估的可能。因此，这个比较暴露的是一组假设之间的不一致。几何映射、质量参数及后来发展的配对或诱导响应修正，需要分别检验；历史对照并未唯一分离它们各自的贡献。值得保留的历史经验是：同一模型同时连接能量和跃迁后，原本能被单独拟合掩盖的限制变得可见。

\subsection{从参考形状继续追问响应}\label{sec:1-3-4}

沿用第~\ref{sec:1-2}~节的定义，实验室静四极矩 $Q_s$ 是具有确定角动量的态在最大磁投影下的电四极对角矩，跃迁约化矩阵元描述不同态之间的电荷连接。这些量约束量子态的电荷关联；平均场极小处的密度是另一层近似参考。集体惯量则进入所选坐标的动能，回答怎样运动需要多少能量；计算它还要考虑扰动引起的密度与场的变化，即残余相互作用参与的诱导响应。静态电荷矩、惯量与响应之间有联系，却不能互相替代。

Hinohara 在 2015 年给出了一种响应计算的例子。在含配对的自洽平均场及相应准粒子随机相位近似（QRPA）中，连续对称性破缺的参考态存在沿对称变换改变的方向。理想自洽响应沿这一方向没有恢复力，产生零频 Nambu--Goldstone（NG）模。它属于参考态及其小振幅响应的描述，并非有限核实验测到的无质量粒子；不完全自洽和基底截断还可使数值伪模偏离零频~\cite{hinohara-2015-nambu-goldstone-inertia-response}。

该文利用有限振幅法，由相应生成元的零频响应求 Thouless--Valatin 惯量，并以配对转动等计算检验。配对转动涉及粒子数规范空间，不等同于核的空间转动。其推进在于把惯量与具体参考态、残余作用和诱导场联系起来；它仍是有近似和输入条件的理论路线，不能据此概括所有集体谱学现象，更不是理论发展的终点。

\begin{table}[htbp]
\centering
\small
\caption{早期统一模型到响应微观理论的核心问题、自由度与适用条件}\label{tab:1-3-unified-to-response}
\begin{tabularx}{\textwidth}{p{3.5cm}p{3.5cm}p{3.5cm}X}
\toprule
问题 & 自由度 & 可计算量 & 条件 \\
\midrule
怎样联系形状与核子运动 & 四极形状、取向与低能粒子态 & 耦合能谱与波函数 & 流动假设、能标与耦合近似 \\
能量与 $E2$ 是否相容 & 同一模型中的转动与电荷算符 & 两条条件形变映射 & 历史采用量及几何、惯量关系 \\
惯量怎样包含微观调整 & 破缺参考态及其诱导响应 & 相应 NG 惯量 & 自洽残余作用、截断与响应通道 \\
\bottomrule
\end{tabularx}
\end{table}

表~\ref{tab:1-3-unified-to-response}所列的问题引出后文的安排：从占据与配对讨论微观结构，再进入形状波函数、惯量和取向模式，并回到实验所能约束的量。第~\ref{sec:1-4}~节先以六个质量区建立比较坐标，说明相同问题在不同核子占据和结构条件下，为什么需要不同的案例与探针。
